\documentclass[pdflatex,sn-mathphys-num]{sn-jnl}

\usepackage{amsmath,amssymb,amsfonts}%
\usepackage{amsthm}%
\usepackage{mathtools}%
\usepackage{enumitem}%
\usepackage{graphicx}%

\hbadness=10000
\vbadness=10000

%% Theorem environments
\theoremstyle{thmstyleone}%
\newtheorem{theorem}{Theorem}[section]
\newtheorem{proposition}[theorem]{Proposition}
\newtheorem{corollary}[theorem]{Corollary}
\newtheorem{lemma}[theorem]{Lemma}
\newtheorem{conjecture}[theorem]{Conjecture}

\theoremstyle{thmstyletwo}%
\newtheorem{remark}[theorem]{Remark}

\theoremstyle{thmstylethree}%
\newtheorem{definition}[theorem]{Definition}

%% Math operators
\DeclareMathOperator{\tr}{tr}
\DeclareMathOperator{\Hess}{Hess}
\DeclareMathOperator{\Ric}{Ric}
\DeclareMathOperator{\OO}{O}
\DeclareMathOperator{\Sym}{Sym}
\DeclareMathOperator{\Hom}{Hom}
\DeclareMathOperator{\Span}{span}

%% Shorthand
\newcommand{\pa}{\partial}
\newcommand{\nn}{\nabla}
\newcommand{\del}{\delta}
\newcommand{\half}{\tfrac{1}{2}}

\numberwithin{equation}{section}

% ================================================================
\begin{document}
% ================================================================

\title[Quaternionic Seam Geometry and Maxwell]{Maxwell's Equations
  from Self-Consistent Quaternionic Seam Geometry}

\author*[1]{\fnm{Lars} \sur{R\"{o}nnb\"{a}ck}}\email{lars@uptochange.com}

\affil*[1]{\orgdiv{Department of Computer and Systems Sciences},
  \orgname{Stockholm University},
  \orgaddress{\city{Stockholm}, \country{Sweden}}}

\abstract{We extend seam geometry to a quaternion-valued
field $\varphi\colon M\to\mathbb{H}$ with $U(1)$-equivariance.
The invariant generators split into five symmetric metric generators
and an antisymmetric 2-form $F = dq_2\wedge dq_3$ that automatically
satisfies the Bianchi identity.  We derive an exact divergence formula
relating $\nabla_\mu F^{\mu\nu}$ to the scalar wave equations and a
Hessian residue, showing that source-free Maxwell's equations hold if
and only if a Hessian orthogonality condition is satisfied.  On a
spherical background, both the azimuthal and polar ans\"{a}tze produce
no-go theorems: the self-consistency condition is incompatible with
non-trivial electromagnetic fields.  On a cylindrical background,
the azimuthal vortex ansatz admits a unique solution $Q = c\rho$,
yielding a uniform magnetic field on a flat self-generated metric.
Restoring the Hessian generator $\gamma_0\nabla^2_h q_0$ breaks this
flatness obstruction: the logarithmic potential $q_0 = d\ln\rho$
produces the exact curved solution
$Q = c(\alpha\rho - \Lambda/\rho)$ with a non-flat metric
$g_{\phi\phi} = (\alpha^2\rho^4-\Lambda^2)/(2\alpha\rho^2)$
that simultaneously satisfies the wave equations and source-free
Maxwell's equations---a self-consistent curved
configuration generated purely from scalars.  The
Einstein tensor of this solution projects onto the
plane orthogonal to~$F$, a relation imposed by
self-consistency rather than postulated.}

\keywords{seam geometry, quaternion, Maxwell equations,
  self-consistency, no-go theorem}

\maketitle

% ================================================================
\section{Introduction}
\label{sec:intro}
% ================================================================

In the companion paper~\cite{seam1} we introduced \emph{seam geometry},
a framework in which a (pseudo-)Riemannian metric on $(M^n,h)$ is
generated from a scalar field $s\colon M\to\mathbb{R}$ through four
$O(n)$-equivariant generators:
\begin{equation}
\label{eq:real-seam}
g = \alpha\,h + \beta\,ds\otimes ds
    + \gamma\,\nn^2 s + \delta\,(\nn^2 s)^2.
\end{equation}
Promoting the scalar to a complex field
$\varphi = u + iv\colon M\to\mathbb{C}$ introduced six generators
and resolved the structural limitations of the real framework, with
the Taub plane-symmetric vacuum providing an explicit non-flat seam
metric.

The natural next step in the hierarchy of division algebras
$\mathbb{R}\subset\mathbb{C}\subset\mathbb{H}$~\cite{BaezHuerta2010} is the
\emph{quaternionic seam}: a field
$\varphi = q_0 + q_1\mathbf{i} + q_2\mathbf{j} + q_3\mathbf{k}
\colon M\to\mathbb{H}$ with four real components.
The general quaternionic seam has 25 equivariant generators at
orders~$\le 2$.  The central question is what internal symmetry
to impose.

Two boundary cases are instructive.  Maximal internal
symmetry---$Sp(1)\cong SU(2)$ acting by left
multiplication---averages all cross-gradient terms, reducing the
quaternionic seam to formally the real seam with $s = |\varphi|$:
only the total gradient $\sum_a dq_a\otimes dq_a$ and the
background $h$ survive as generators.  The fully general case
is computationally intractable.

The intermediate case, $U(1)\subset Sp(1)$, is the
\emph{Goldilocks} equivariance.  Right multiplication by
$e^{\mathbf{i}\theta}$ rotates $(q_2,q_3)$ while fixing
$(q_0,q_1)$, producing exactly 6 generators---the same count as
the complex seam---plus an invariant antisymmetric 2-form
$F = dq_2\wedge dq_3$ that has no analogue in the complex
framework.  This $F$ automatically satisfies the Bianchi identity
$dF = 0$ and is a natural candidate for the electromagnetic field
tensor.  The programme of deriving electromagnetism from
geometric self-consistency has a long history, from
Kaluza--Klein theory~\cite{Kaluza1921,Klein1926} to the
Rainich--Misner--Wheeler ``already unified field
theory''~\cite{Rainich1925,MisnerWheeler1957}.  The seam
approach is complementary: $F$ arises not from an extra
dimension or from algebraic conditions on the Ricci tensor,
but from the internal $U(1)$ structure of a quaternionic
scalar field.

\medskip\noindent
\textbf{Summary of results.}
\begin{enumerate}[noitemsep]
\item An exact divergence formula
  (Theorem~\ref{thm:maxwell-from-seam}) relating
  $\nn_\mu F^{\mu\nu}$ to the scalar wave equations and a
  Hessian residue: Maxwell holds iff a Hessian orthogonality
  condition is satisfied.
\item A no-go theorem (Theorem~\ref{thm:no-go-maxwell}) for the
  azimuthal ansatz on a spherical background: the numerator of
  $\nn_\mu F^{\mu\phi}$ is a sum of non-negative terms that
  forces $F = 0$.
\item A stronger no-go (Theorem~\ref{thm:polar-nogo}) for the
  polar ansatz: the dipole eigenvalue $\ell(\ell+1) = 2$
  introduces an unconditional obstruction.
\item A unique flat escape (Theorem~\ref{thm:cyl-maxwell}): the
  cylindrical vortex ansatz with $Q = c\rho$ produces a
  self-consistent uniform magnetic field on a flat metric.
\item A curved solution (Theorem~\ref{thm:curved-cyl}): restoring
  the Hessian generator yields the exact solution
  $Q = c(\alpha\rho - \Lambda/\rho)$ on a non-flat metric with
  a curvature singularity at $\rho_0 = (\Lambda/\alpha)^{1/2}$.
\item A curvature analysis (Proposition~\ref{prop:curvature}):
  the Einstein tensor satisfies
  $F_{\mu\alpha}\,G^{\alpha\nu} = 0$---it vanishes in the
  electromagnetic plane and projects onto the orthogonal
  $(t,z)$-subspace.  Standard Einstein--Maxwell is
  \emph{not} satisfied; instead, self-consistency imposes an
  emergent projection identity.
\end{enumerate}

% ================================================================
\section{The \texorpdfstring{$U(1)$}{U(1)}-equivariant quaternionic seam}
\label{sec:u1-seam}
% ================================================================

\subsection{\texorpdfstring{$U(1)$}{U(1)} action and invariant generators}

The $U(1)$ subgroup acts by right multiplication by
$e^{\mathbf{i}\theta}$, i.e.\
$(q_0,q_1,q_2,q_3)\mapsto (q_0,q_1,
  q_2\cos\theta - q_3\sin\theta,\,
  q_2\sin\theta + q_3\cos\theta)$.
The fields split into an \emph{invariant pair} $(q_0,q_1)$ and a
\emph{rotating pair} $(q_2,q_3)$.

\begin{proposition}[$U(1)$-invariant gradient generators]
\label{prop:u1-generators}
The $U(1)$-invariant symmetric $(0,2)$-tensors built from the first
derivatives of $(q_0,q_1,q_2,q_3)$ are:
\begin{enumerate}[noitemsep,label=(\roman*)]
\item $G_{00} = \pa_\mu q_0\,\pa_\nu q_0$
  \hfill (from invariant pair)
\item $G_{11} = \pa_\mu q_1\,\pa_\nu q_1$
  \hfill (from invariant pair)
\item $G_{01} = \pa_{(\mu}q_0\,\pa_{\nu)}q_1$
  \hfill (cross-term, invariant pair)
\item $G_\perp = \pa_\mu q_2\,\pa_\nu q_2 + \pa_\mu q_3\,\pa_\nu q_3$
  \hfill ($U(1)$-invariant from rotating pair)
\item $h_{\mu\nu}$ \hfill (conformal generator)
\end{enumerate}
No cross-gradients between the invariant pair and the rotating pair
survive: $\pa_{(\mu}q_0\,\pa_{\nu)}q_2$ transforms under $U(1)$ as
the first component of a doublet and averages to zero.
The total count is \textbf{5 gradient generators $+$ 1 conformal
$=$ 6 generators}---the same as the complex seam.
\end{proposition}

\begin{proof}
Under $\theta$-rotation, $\pa_\mu q_2\mapsto \cos\theta\,\pa_\mu q_2
- \sin\theta\,\pa_\mu q_3$ and
$\pa_\mu q_3\mapsto\sin\theta\,\pa_\mu q_2
+ \cos\theta\,\pa_\mu q_3$.  Direct computation shows
$\pa_\mu q_2\,\pa_\nu q_2 + \pa_\mu q_3\,\pa_\nu q_3$ is invariant
(the rotation is orthogonal).  The antisymmetric combination
$\pa_\mu q_2\,\pa_\nu q_3 - \pa_\mu q_3\,\pa_\nu q_2$ is also
invariant but is not a metric generator (it is the 2-form $F$;
see below).  All other bilinears involving exactly one rotated
index transform non-trivially.
\end{proof}

\subsection{The invariant 2-form}

Besides the metric generators, the $U(1)$-invariant 2-form
\begin{equation}
\label{eq:u1-Fform}
F_{\mu\nu} := \pa_\mu q_2\,\pa_\nu q_3 - \pa_\nu q_2\,\pa_\mu q_3
            = (dq_2\wedge dq_3)_{\mu\nu}
\end{equation}
is the antisymmetric counterpart of the surviving generators.

\begin{proposition}[Automatic Bianchi identity]
\label{prop:bianchi}
The 2-form $F = dq_2\wedge dq_3$ satisfies $dF = 0$ identically.
\end{proposition}

\begin{proof}
$dF = d(dq_2\wedge dq_3) = d^2q_2\wedge dq_3 - dq_2\wedge d^2q_3
= 0$, since $d^2 = 0$ on any smooth function.
\end{proof}

\noindent
Thus $F$ automatically satisfies the homogeneous Maxwell equation.
Whether the inhomogeneous equation $\nn_\mu F^{\mu\nu} = 0$ holds
depends on the metric and the field equations for $q_2,q_3$.

\subsection{The seam metric}

\begin{definition}[$U(1)$-equivariant quaternionic seam metric]
\label{def:u1-seam-metric}
The $U(1)$-equivariant quaternionic seam metric of order~$\le 1$
(gradient generators only) is
\begin{equation}
\label{eq:u1-seam-metric}
g_{\mu\nu}
= \alpha\,h_{\mu\nu}
  + \beta_0\,\pa_\mu q_0\,\pa_\nu q_0
  + \beta_1\,\pa_\mu q_1\,\pa_\nu q_1
  + \beta_\times\,\pa_{(\mu}q_0\,\pa_{\nu)}q_1
  + \beta_\perp\bigl(
      \pa_\mu q_2\,\pa_\nu q_2 + \pa_\mu q_3\,\pa_\nu q_3
    \bigr),
\end{equation}
where the coefficients $\alpha,\beta_0,\beta_1,\beta_\times,
\beta_\perp$ are smooth functions of the $U(1)$-invariant
scalar jet data.
\end{definition}

\begin{remark}[Comparison with the complex seam]
\label{rem:u1-vs-complex}
The complex seam of~\cite{seam1} has six generators built from
two real fields $u,v$.  The $U(1)$-equivariant quaternionic seam
has the same count with a different origin: the key structural
difference is that the antisymmetric 2-form $F = dq_2\wedge dq_3$
has no analogue in the two-real-field complex seam.
\end{remark}

% ================================================================
\section{Self-consistency and Maxwell's equations}
\label{sec:maxwell}
% ================================================================

The seam metric is generated by scalar fields, so there is no
\emph{a~priori} guarantee that the resulting geometry supports
electromagnetic propagation.  We now ask: under what conditions
does the wave equation $\square_g q_a = 0$ imply the source-free
Maxwell equation $\nn_\mu F^{\mu\nu} = 0$?  The answer is an
exact divergence identity (Theorem~\ref{thm:maxwell-from-seam})
that reduces Maxwell to a single Hessian alignment condition.

\subsection{The divergence identity}

\begin{definition}[$U(1)$-seam self-consistency]
\label{def:u1-self-consistency}
The $U(1)$-equivariant quaternionic seam is \emph{self-consistent}
if each scalar field satisfies its wave equation on the metric it
generates:
\begin{equation}
\label{eq:u1-self-consistency}
\square_g q_a
:= \frac{1}{\sqrt{|g|}}\,\pa_\mu\bigl(\sqrt{|g|}\,g^{\mu\nu}\,
   \pa_\nu q_a\bigr) = m_a^2\,q_a,
\qquad a=0,1,2,3,
\end{equation}
where the $m_a$ are mass parameters
(with $m_a = 0$ for the massless case).
\end{definition}

\begin{theorem}[Divergence of $F$ from scalar wave equations]
\label{thm:maxwell-from-seam}
Let $(M,g)$ be a pseudo-Riemannian manifold and
$q_2,q_3\in C^\infty(M)$.  Define $F := dq_2\wedge dq_3$.
Then
\begin{equation}
\label{eq:div-F-exact}
\nn_\mu F^{\mu\nu}
= (\square_g q_2)\,\pa^\nu q_3
  - (\square_g q_3)\,\pa^\nu q_2
  - \bigl[
      (\Hess\,q_2)^{\nu\mu}\,\pa_\mu q_3
    - (\Hess\,q_3)^{\nu\mu}\,\pa_\mu q_2
    \bigr],
\end{equation}
where $(\Hess\,q_a)^{\nu\mu} := g^{\nu\beta}\nn_\beta\nn^\mu q_a$.

In particular, $F$ satisfies $\nn_\mu F^{\mu\nu} = 0$ if and only if
the \textbf{Hessian orthogonality condition}
\begin{equation}
\label{eq:hessian-orthogonality}
(\square_g q_2)\,\pa^\nu q_3
  - (\square_g q_3)\,\pa^\nu q_2
= (\Hess\,q_2)^{\nu\mu}\,\pa_\mu q_3
  - (\Hess\,q_3)^{\nu\mu}\,\pa_\mu q_2
\end{equation}
holds.
When both fields are massless ($\square_g q_a = 0$), this reduces to
\begin{equation}
\label{eq:hessian-condition-massless}
(\Hess\,q_2)^{\nu\mu}\,\pa_\mu q_3
= (\Hess\,q_3)^{\nu\mu}\,\pa_\mu q_2,
\end{equation}
which is \textbf{not} automatically satisfied.
\end{theorem}

\begin{proof}
The Bianchi identity $dF = 0$ is automatic
(Proposition~\ref{prop:bianchi}).  For the divergence, compute:
\begin{align}
\nn_\mu F^{\mu\nu}
&= \nn_\mu\bigl(g^{\mu\alpha}g^{\nu\beta}
   (\pa_\alpha q_2\,\pa_\beta q_3 - \pa_\alpha q_3\,\pa_\beta q_2)
   \bigr).
  \label{eq:div-F-expand}
\end{align}
Using the Leibniz rule and $\nn_\mu(g^{\mu\alpha}\pa_\alpha q_a)
= \square_g q_a$:
\begin{align}
\nn_\mu F^{\mu\nu}
&= g^{\nu\beta}\bigl[
     (\square_g q_2)\,\pa_\beta q_3
   + (\nn q_2)^\mu\,\nn_\mu\pa_\beta q_3
   - (\square_g q_3)\,\pa_\beta q_2
   - (\nn q_3)^\mu\,\nn_\mu\pa_\beta q_2
   \bigr].
  \label{eq:div-F-leibniz}
\end{align}
Since $\nn_\mu\pa_\beta q_a = \nn_\beta\pa_\mu q_a$ (torsion-free
connection), the gradient--Hessian terms become:
\begin{align*}
(\nn q_2)^\mu\nn_\beta\pa_\mu q_3
&= \nn_\beta\bigl(\langle\nn q_2,\nn q_3\rangle_g\bigr)
   - (\nn_\beta\nn^\mu q_2)\,\pa_\mu q_3, \\
(\nn q_3)^\mu\nn_\beta\pa_\mu q_2
&= \nn_\beta\bigl(\langle\nn q_3,\nn q_2\rangle_g\bigr)
   - (\nn_\beta\nn^\mu q_3)\,\pa_\mu q_2.
\end{align*}
The $\nn_\beta\langle\cdots\rangle$ terms cancel (the inner product
is symmetric).  The Hessian residue survives,
giving~\eqref{eq:div-F-exact}.
\end{proof}

\begin{remark}[Counterexample: massless waves do not suffice]
\label{rem:counterexample}
On flat Minkowski space, $q_2 = \sin(t+x)$ and
$q_3 = \sin(t+y)$ both satisfy $\square q_a = 0$, yet
$\nn_\mu F^{\mu t} = -\cos(t+x)\sin(t+y) \neq 0$.
The Hessians point in different null directions, violating
the alignment condition~\eqref{eq:hessian-condition-massless}.
\end{remark}

\subsection{The Hessian condition on a static spherical ansatz}

We specialise to a static ansatz on the $n=4$ Minkowski background
with coordinates $(t,r,\theta,\phi)$:
\begin{equation}
\label{eq:u1-ansatz}
q_0 = q_0(r), \qquad q_1 = \omega\,t, \qquad
q_2 = Q(r)\cos\phi, \qquad q_3 = Q(r)\sin\phi,
\end{equation}
where $\omega$ is a constant and $Q(r)$ is a radial profile.
The $\phi$-winding gives a $U(1)$-equivariant configuration with
\begin{equation}
\label{eq:F-ansatz}
F = dq_2\wedge dq_3 = Q\,Q'\,dr\wedge d\phi.
\end{equation}

The $U(1)$-seam metric~\eqref{eq:u1-seam-metric} evaluates to a
diagonal metric (the cross-term $G_{01} = 0$ since the gradients
of $q_0$ and $q_1$ are orthogonal):
\begin{equation}
\label{eq:u1-metric-explicit}
g = \operatorname{diag}\!\bigl(
  -\alpha + \beta_1\omega^2,\;\;
  \alpha + \beta_0(q_0')^2 + \beta_\perp(Q')^2,\;\;
  \alpha\,r^2,\;\;
  \alpha\,r^2\sin^2\!\theta + \beta_\perp Q^2
\bigr).
\end{equation}

\begin{proposition}[Hessian condition on the ansatz]
\label{prop:hessian-ansatz}
For the ansatz $(q_2,q_3) = Q(r)(\cos\phi,\sin\phi)$ on a
diagonal metric, the Hessian orthogonality
condition~\eqref{eq:hessian-condition-massless} satisfies:
\begin{enumerate}[noitemsep,label=(\roman*)]
\item For $\nu\in\{t,r,\theta\}$, the condition holds identically
  (by $U(1)$ averaging: $\sin\phi\cos\phi$ cancels).
\item For $\nu = \phi$, the condition is a single radial ODE
  relating $Q(r)$, $Q'(r)$, and the metric components
  $g_{rr}$, $g_{\phi\phi}$.
\end{enumerate}
\end{proposition}

\begin{proof}
The gradients $\pa_\mu q_{2,3}$ have support only in the $r$
and $\phi$ directions.  The diagonal Hessian components
$H_{rr}(q_2) = Q''\cos\phi$ and $H_{rr}(q_3) = Q''\sin\phi$
contribute terms proportional to
$\cos\phi\sin\phi - \sin\phi\cos\phi = 0$ by $U(1)$ rotation,
so the $\nu\in\{t,r,\theta\}$ components vanish.  The
off-diagonal Hessians $H_{r\phi}(q_2) = -(Q'-Q/r)\sin\phi$ and
$H_{r\phi}(q_3) = (Q'-Q/r)\cos\phi$ contribute only to
$\nu = \phi$ after $\phi$-averaging ($\sin^2\phi + \cos^2\phi = 1$).
\end{proof}

\subsection{Electromagnetic stress-energy}

The 2-form $F = Q\,Q'\,dr\wedge d\phi$ produces a diagonal
electromagnetic stress-energy tensor
\begin{equation}
\label{eq:T-em-explicit}
T^{(F)}_{\mu\nu} = F_{\mu\alpha}F_\nu{}^\alpha
  - \tfrac{1}{4}\,g_{\mu\nu}\,F_{\alpha\beta}F^{\alpha\beta}.
\end{equation}

\begin{proposition}[$T^{(F)}$ is traceless]
\label{prop:T-traceless}
The stress-energy tensor~\eqref{eq:T-em-explicit} satisfies
$g^{\mu\nu}T^{(F)}_{\mu\nu} = 0$ in $n=4$ dimensions.
\end{proposition}

\begin{proof}
The trace $F_{\mu\alpha}F^{\mu\alpha}
= F_{\alpha\beta}F^{\alpha\beta}$ cancels the
$\frac{1}{4}\cdot 4 = 1$ times $F^2$ from the second term.
\end{proof}

% ================================================================
\section{No-go theorems on spherical backgrounds}
\label{sec:nogo}
% ================================================================

With the divergence identity in hand, we now test whether
specific field configurations can satisfy both the wave equations
and Maxwell simultaneously.  The most natural starting point is
a spherically symmetric background, where the rotating pair
inherits the angular structure of the coordinate system.

\subsection{No-go for the azimuthal ansatz}

The azimuthal ansatz~\eqref{eq:u1-ansatz}, with $q_{2,3}$
winding in the $\phi$-direction, is the simplest
$U(1)$-equivariant configuration on a spherical background.

\begin{theorem}[No-go: azimuthal ansatz on spherical background]
\label{thm:no-go-maxwell}
Let $g$ be the $U(1)$-equivariant seam
metric~\eqref{eq:u1-metric-explicit} with $\alpha > 0$ and
$\beta_0 \ge 0$, generated by the
ansatz~\eqref{eq:u1-ansatz}.  If $q_2$ and $q_3$ satisfy the
massless wave equation $\square_g q_a = 0$, then
\begin{equation}
\label{eq:no-go-residue}
\nn_\mu F^{\mu\phi}
= \frac{
    \alpha\,Q^2\cos^2\!\theta
  + \beta_0\,Q^2(q_0')^2
  + \alpha\,\sin^2\!\theta\,(rQ'-Q)^2
  }{g_{rr}\,g_{\phi\phi}^2}.
\end{equation}
Each term is non-negative, so $\nn_\mu F^{\mu\nu} = 0$ implies
$Q \equiv 0$ and hence $F = 0$.  In particular, the
$U(1)$-equivariant seam with this ansatz admits \textbf{no
non-trivial self-consistent Maxwell solution on a curved seam metric}.
\end{theorem}

\begin{proof}
\emph{Step~1: Maxwell divergence via the wave equation.}
The only non-zero component of $F^{\mu\nu}$ is
$F^{r\phi} = QQ'/(g_{rr}\,g_{\phi\phi})$, so
\begin{equation}
\label{eq:div-F-explicit}
\nn_\mu F^{\mu\phi}
= \frac{1}{\sqrt{|g|}}\,\pa_r\!\left(
    \sqrt{|g|}\,\frac{QQ'}{g_{rr}\,g_{\phi\phi}}
  \right).
\end{equation}
By the product rule,
\begin{equation}
\label{eq:maxwell-product-rule}
\nn_\mu F^{\mu\phi}
= \frac{Q}{g_{\phi\phi}}\cdot
  \frac{1}{\sqrt{|g|}}\,\pa_r\!\left(
    \frac{\sqrt{|g|}\,Q'}{g_{rr}}
  \right)
\;+\; \frac{Q'}{g_{rr}}\,\pa_r\!\left(
    \frac{Q}{g_{\phi\phi}}
  \right).
\end{equation}
The wave equation $\square_g(Q\cos\phi) = 0$ on the diagonal metric
gives (separating $\cos\phi$, with $-Q/g_{\phi\phi}$ from
$\pa_\phi^2\cos\phi = -\cos\phi$):
\begin{equation}
\label{eq:wave-Q}
\frac{1}{\sqrt{|g|}}\,\pa_r\!\left(
  \frac{\sqrt{|g|}\,Q'}{g_{rr}}\right)
= \frac{Q}{g_{\phi\phi}}.
\end{equation}
Substituting:
\begin{equation}
\label{eq:maxwell-after-wave}
\nn_\mu F^{\mu\phi}
= \frac{Q^2\,g_{rr} + (Q')^2\,g_{\phi\phi}
        - Q'Q\,g_{\phi\phi}'}
       {g_{rr}\,g_{\phi\phi}^2}.
\end{equation}

\medskip
\emph{Step~2: Expansion of the numerator.}
With $g_{rr} = \alpha + \beta_0(q_0')^2 + \beta_\perp(Q')^2$,
$g_{\phi\phi} = \alpha r^2\sin^2\!\theta + \beta_\perp Q^2$,
and $g_{\phi\phi}' = 2\alpha r\sin^2\!\theta + 2\beta_\perp QQ'$:
\begin{align*}
\mathcal{N}
&:= Q^2\,g_{rr} + (Q')^2\,g_{\phi\phi} - Q'Q\,g_{\phi\phi}' \\
&= \alpha Q^2 + \beta_0 Q^2(q_0')^2
   + \alpha\sin^2\!\theta\,
     (r^2(Q')^2 - 2rQ'Q).
\end{align*}
(The $\beta_\perp$ terms cancel exactly.)
Completing the square:
$r^2(Q')^2 - 2rQ'Q = (rQ'-Q)^2 - Q^2$, so
\begin{equation*}
\mathcal{N}
= \alpha Q^2\cos^2\!\theta + \beta_0 Q^2(q_0')^2
  + \alpha\sin^2\!\theta\,(rQ'-Q)^2.
\end{equation*}

\medskip
\emph{Step~3: Positivity.}
With $\alpha > 0$ and $\beta_0 \ge 0$, all three terms are
$\ge 0$.  Since $\cos^2\!\theta$ and $1$ are linearly independent
on $(0,\pi)$, both coefficients must vanish separately.
The constant coefficient is a sum of non-negative terms, so
$Q(r) = 0$ for all $r$, giving $F = 0$.
\end{proof}

\begin{remark}[Source of the obstruction]
\label{rem:obstruction-source}
The obstruction is geometric: the Minkowski background in
spherical coordinates has $g_{\phi\phi}^{(\eta)}
= r^2\sin^2\!\theta$, but the seam correction
$\beta_\perp Q^2$ is $\theta$-independent.  The mismatch
produces a $\cos^2\!\theta$ term in $\mathcal{N}$ that no
radial profile $Q(r)$ can cancel.
\end{remark}

\subsection{No-go for the polar ansatz}

\begin{definition}[Polar ansatz]
\label{def:polar-ansatz}
On Minkowski spacetime in spherical coordinates, set
\begin{equation}
\label{eq:polar-ansatz}
q_0 = q_0(r), \qquad q_1 = \omega\,t, \qquad
q_2 = Q(r)\cos\theta, \qquad q_3 = Q(r)\sin\theta.
\end{equation}
\end{definition}

The seam metric becomes
\begin{equation}
\label{eq:polar-metric}
g = \operatorname{diag}\!\bigl(
  -\alpha+\beta_1\omega^2,\;\;
  \alpha + \beta_0(q_0')^2 + \beta_\perp(Q')^2,\;\;
  \alpha r^2 + \beta_\perp Q^2,\;\;
  \alpha r^2\sin^2\!\theta
\bigr),
\end{equation}
with $F = QQ'\,dr\wedge d\theta$.

\begin{theorem}[Stronger no-go: polar ansatz]
\label{thm:polar-nogo}
On the polar ansatz~\eqref{eq:polar-ansatz} with $\alpha > 0$,
$\beta_0\ge 0$, $\beta_\perp\ge 0$, and $\square_g q_{2,3} = 0$,
\begin{equation}
\label{eq:polar-numerator}
\nn_\mu F^{\mu\theta}
= \frac{\alpha\bigl[Q^2 + (rQ'-Q)^2\bigr]
        + 2\beta_0\,Q^2(q_0')^2
        + \beta_\perp\,Q^2(Q')^2}
       {g_{rr}\,g_{\theta\theta}^2}.
\end{equation}
The numerator is strictly positive for $Q\neq 0$, so
$Q\equiv 0$ is forced.
\end{theorem}

\begin{proof}
The wave equation $\square_g(Q\cos\theta) = 0$ gives
\begin{equation}
\label{eq:wave-polar}
\frac{1}{\sqrt{|g|}}\,\pa_r\!\left(
  \frac{\sqrt{|g|}\,Q'}{g_{rr}}\right)
= \frac{2Q}{g_{\theta\theta}},
\end{equation}
where the factor $2$ (vs.\ $1$ in~\eqref{eq:wave-Q}) arises
because $\cos\theta$ is a dipole harmonic: eigenvalue
$\ell(\ell+1) = 2$ instead of $m^2 = 1$ for $\cos\phi$.

The product-rule decomposition gives numerator
\begin{align*}
\mathcal{N}
&= 2Q^2\,g_{rr} + (Q')^2\,g_{\theta\theta}
   - QQ'\,g_{\theta\theta}' \\
&= \alpha\bigl[Q^2 + (rQ'-Q)^2\bigr]
 + 2\beta_0 Q^2(q_0')^2
 + \beta_\perp Q^2(Q')^2.
\end{align*}
Since $Q^2 + (rQ'-Q)^2 \ge Q^2 > 0$ for $Q\neq 0$, and all
other terms are non-negative, $\mathcal{N} = 0$ requires $Q = 0$.
\end{proof}

\begin{remark}[Eigenvalue mismatch]
\label{rem:polar-stronger}
The azimuthal ansatz had eigenvalue $m^2 = 1$ from $\cos\phi$,
allowing the cylindrical background
(Section~\ref{sec:cyl}) to zero the numerator with
$Q = c\rho$.  The polar eigenvalue $\ell(\ell+1) = 2$ produces
an extra $\alpha Q^2$ term that cannot be absorbed by any radial
profile, making the no-go unconditional.
\end{remark}

\subsection{The hedgehog obstruction}

\begin{remark}[Hedgehog ansatz]
\label{rem:hedgehog}
A \emph{hedgehog ansatz} $q_a = f(r)\,x_a/r$ ($a=1,2,3$)
maps each spatial direction to an imaginary quaternion component.
However, a spherically symmetric monopole field requires a
\emph{non-abelian} field strength $F^a_{\mu\nu}$ with all three
internal components~\cite{tHooft1974,Polyakov1974}; the abelian
$F = dq_2\wedge dq_3$ selects a plane in internal space and
breaks spherical symmetry.
Since $Sp(1)$-equivariance collapses to the real seam (only
$\sum_a dq_a\otimes dq_a$ survives), the hedgehog route would
require an intermediate equivariance group strictly between
$U(1)$ and $Sp(1)$.
\end{remark}

% ================================================================
\section{The cylindrical vortex solution}
\label{sec:cyl}
% ================================================================

The spherical no-go theorems trace to a common source: the
background angular structure ($\sin^2\!\theta$ in the azimuthal
case, $\ell(\ell+1)=2$ in the polar case) produces positive-definite
obstructions that no radial profile can cancel.  A cylindrical
background eliminates this angular mismatch, opening the door to
non-trivial solutions.

\begin{definition}[Cylindrical $U(1)$-seam ansatz]
\label{def:cylindrical-ansatz}
On Minkowski spacetime in cylindrical coordinates
$(t,\rho,\phi,z)$ with $\eta = \operatorname{diag}(-1,1,\rho^2,1)$,
set
\begin{equation}
\label{eq:cyl-ansatz}
q_0 = q_0(\rho), \qquad q_1 = \omega\,t, \qquad
q_2 = Q(\rho)\cos\phi, \qquad q_3 = Q(\rho)\sin\phi.
\end{equation}
\end{definition}

The $U(1)$-seam metric~\eqref{eq:u1-seam-metric} evaluates to
\begin{equation}
\label{eq:cyl-metric}
g = \operatorname{diag}\!\bigl(
  -\alpha+\beta_1\omega^2,\;\;
  \alpha+\beta_0(q_0')^2+\beta_\perp(Q')^2,\;\;
  \alpha\rho^2+\beta_\perp Q^2,\;\;
  \alpha
\bigr),
\end{equation}
where primes denote $d/d\rho$.  Crucially, $g_{\phi\phi}
= \alpha\rho^2 + \beta_\perp Q^2$ contains no $\sin^2\!\theta$
factor---the angular mismatch of
Remark~\ref{rem:obstruction-source} is absent.

\begin{theorem}[Cylindrical Maxwell from self-consistency]
\label{thm:cyl-maxwell}
On the cylindrical ansatz~\eqref{eq:cyl-ansatz}, imposing the wave
equation $\square_g q_{2,3} = 0$ gives
\begin{equation}
\label{eq:cyl-numerator}
\nn_\mu F^{\mu\phi}
= \frac{\alpha(\rho\,Q' - Q)^2
        + \beta_0\,Q^2(q_0')^2}
       {g_{\rho\rho}\,g_{\phi\phi}^2}.
\end{equation}
This vanishes for all $\rho > 0$ if and only if
\begin{equation}
\label{eq:cyl-maxwell-conditions}
Q(\rho) = c\,\rho
\qquad\text{and}\qquad
q_0'(\rho) = 0
\quad(\text{or } Q \equiv 0).
\end{equation}
\end{theorem}

\begin{proof}
The computation follows the same product-rule strategy as
Theorem~\ref{thm:no-go-maxwell} with $r\to\rho$ and no
$\theta$-dependence.  The numerator is:
\begin{align*}
\mathcal{N}
&= Q^2\,g_{\rho\rho} + (Q')^2\,g_{\phi\phi} - QQ'\,g_{\phi\phi}'
\\
&= \alpha(\rho Q' - Q)^2 + \beta_0 Q^2(q_0')^2,
\end{align*}
where the $\beta_\perp$ terms cancel exactly.  With $\alpha > 0$
and $\beta_0 \ge 0$, both terms are non-negative.  Setting
$\mathcal{N} = 0$ requires $\rho Q' = Q$ (i.e.\ $Q = c\rho$) and
either $q_0' = 0$ or $Q = 0$.
\end{proof}

\begin{proposition}[Self-consistent cylindrical solution]
\label{prop:cyl-solution}
The configuration $Q = c\rho$, $q_0 = \mathrm{const}$,
$q_1 = \omega t$ satisfies simultaneously:
\begin{enumerate}[noitemsep,label=(\roman*)]
\item the wave equations $\square_g q_a = 0$ for $a = 0,1,2,3$;
\item the source-free Maxwell equations $\nn_\mu F^{\mu\nu} = 0$
  and $dF = 0$;
\item the seam metric is
\begin{equation}
\label{eq:cyl-solution-metric}
g = \operatorname{diag}\!\bigl(
  -\alpha+\beta_1\omega^2,\;\;
  A,\;\;
  A\,\rho^2,\;\;
  \alpha
\bigr),
\qquad
A := \alpha + \beta_\perp c^2;
\end{equation}
\item the electromagnetic field is
\begin{equation}
\label{eq:cyl-F}
F = c^2\rho\,d\rho\wedge d\phi
  = c^2\,dx\wedge dy,
\end{equation}
a uniform magnetic field along the $z$-axis.
\end{enumerate}
\end{proposition}

\begin{proof}
(i) $q_0 = \mathrm{const}$: $\square_g q_0 = 0$ trivially.
$q_1 = \omega t$: $\square_g q_1
= \omega\,\pa_t(\sqrt{|g|}\,g^{tt})/\sqrt{|g|} = 0$ since $g^{tt}$
and $\sqrt{|g|}$ are $t$-independent.
For $q_2 = c\rho\cos\phi$: the wave equation reduces
to~\eqref{eq:wave-Q} (with $r\to\rho$).  Setting
$Q' = c$, $g_{\rho\rho} = A$, $g_{\phi\phi} = A\rho^2$:
\[
\text{LHS} = \frac{1}{\rho K}\pa_\rho\!\left(\frac{\rho K c}{A}\right)
= \frac{c}{A\rho},
\qquad
\text{RHS} = \frac{c\rho}{A\rho^2} = \frac{c}{A\rho},
\]
where $K = \sqrt{|g_{tt}|\,A^2\,\alpha}$ is $\rho$-independent.
The equation for $q_3$ is identical by $U(1)$ rotation.

(ii) Theorem~\ref{thm:cyl-maxwell} + Proposition~\ref{prop:bianchi}.

(iii) Substitute $Q = c\rho$, $q_0' = 0$
into~\eqref{eq:cyl-metric}.

(iv) $F = QQ'\,d\rho\wedge d\phi = c^2\rho\,d\rho\wedge d\phi
= c^2\,dx\wedge dy$.
\end{proof}

\begin{remark}[Physical interpretation]
\label{rem:cyl-interpretation}
The solution describes a uniform magnetic field
$\mathbf{B} = c^2\,\hat{z}$ on a static spacetime with
anisotropic spatial scaling: the $(x,y)$-plane is scaled by
$A = \alpha+\beta_\perp c^2$ while the $z$-direction is scaled
by $\alpha$.  The metric~\eqref{eq:cyl-solution-metric} has
constant coefficients, so $G_{\mu\nu} = 0$: the electromagnetic
stress-energy does \emph{not} curve spacetime in this solution.
Whether non-constant coefficient functions $\alpha(Q^2)$,
$\beta_\perp(Q^2)$ produce gravitational backreaction (a
Melvin-type magnetic universe~\cite{Melvin1964}) remains open.
\end{remark}

\begin{remark}[Uniqueness]
\label{rem:cyl-uniqueness}
Within the class of constant seam coefficients,
$Q = c\rho$ is the \emph{unique} non-trivial radial profile
satisfying both the wave equation and Maxwell's equations on
the self-generated seam metric.  The coupled system is exactly
determined: the wave equation fixes $Q$, and the Maxwell
condition further requires $q_0 = \mathrm{const}$.
\end{remark}

% ================================================================
\section{Restoring Hessian generators: curvature from
  self-consistency}
\label{sec:hessian}
% ================================================================

The cylindrical solution
(Proposition~\ref{prop:cyl-solution}) has a flat metric
($G_{\mu\nu} = 0$) despite a non-zero electromagnetic field.
In the companion paper~\cite{seam1}, curvature with constant
seam coefficients arose from the
\emph{Hessian generators} $\gamma\,\nn^2_h s$ and
$\delta\,(\nn^2_h s)^2$: the Hessian of a radial field $s(r)$
is position-dependent even when $\gamma,\,\delta$ are constants
(e.g.\ $(\nn^2_h s)_{\theta\theta} = r\,s'$ on the spherical
background).  The $U(1)$-equivariant metric of
Definition~\ref{def:u1-seam-metric} was restricted to
order~$\le 1$ (gradient generators only), deliberately
excluding the mechanism that produced curvature in the real
seam.

\subsection{\texorpdfstring{$U(1)$}{U(1)}-invariant Hessian generators}

The $U(1)$-invariant Hessian generators at order~$2$ include:
\begin{enumerate}[noitemsep,label=(\roman*)]
\item $(\nn^2_h q_0)_{\mu\nu}$ --- Hessian of the invariant
  field $q_0$ (individually $U(1)$-invariant),
\item $(\nn^2_h q_1)_{\mu\nu}$ --- Hessian of $q_1$
  (individually $U(1)$-invariant),
\item $(\nn^2_h q_2)_{\mu\alpha}\,(\nn^2_h q_2)^\alpha{}_\nu
  + (\nn^2_h q_3)_{\mu\alpha}\,(\nn^2_h q_3)^\alpha{}_\nu$
  --- $U(1)$-averaged squared Hessian from the rotating pair.
\end{enumerate}
The individual Hessians $(\nn^2_h q_2)_{\mu\nu}$ and
$(\nn^2_h q_3)_{\mu\nu}$ transform as a doublet under $U(1)$
and cannot appear separately (just as for the gradients), but
generators~(i) and~(ii) are available directly.

For the cylindrical ansatz with $q_1 = \omega t$, the Hessian
$\nn^2_h(\omega t) = 0$ (all Christoffel symbols
$\Gamma^t_{\mu\nu} = 0$ on the flat background), so
generator~(ii) vanishes.  The simplest non-trivial extension is
therefore to include the Hessian of $q_0$:

\begin{definition}[Extended $U(1)$-seam metric with Hessian
  generator]
\label{def:u1-seam-hessian}
The extended $U(1)$-equivariant quaternionic seam metric of
order~$\le 2$ is
\begin{equation}
\label{eq:u1-seam-hessian}
g_{\mu\nu}
= \alpha\,h_{\mu\nu}
  + \beta_0\,\pa_\mu q_0\,\pa_\nu q_0
  + \beta_1\,\pa_\mu q_1\,\pa_\nu q_1
  + \beta_\perp\,G_{\perp\,\mu\nu}
  + \gamma_0\,(\nn^2_h q_0)_{\mu\nu},
\end{equation}
where $G_\perp$, $\alpha$, $\beta_0$, $\beta_1$, $\beta_\perp$,
$\gamma_0$ are as before (the cross-term $\beta_\times$ drops
out on our ansatz).
\end{definition}

\subsection{Cylindrical metric with the Hessian generator}

On the cylindrical background $h = \operatorname{diag}(-1,1,\rho^2,1)$,
the only non-zero components of $\nn^2_h q_0$ for
$q_0 = q_0(\rho)$ are:
\begin{equation}
\label{eq:hess-q0-cyl}
(\nn^2_h q_0)_{\rho\rho} = q_0'',
\qquad
(\nn^2_h q_0)_{\phi\phi}
= -\Gamma^\rho_{\phi\phi}\,q_0' = \rho\,q_0',
\end{equation}
where $\Gamma^\rho_{\phi\phi} = -\rho$ is the background
Christoffel symbol.  The extended metric is therefore
\begin{equation}
\label{eq:cyl-metric-hess}
g = \operatorname{diag}\!\bigl(
  -\alpha+\beta_1\omega^2,\;\;
  \alpha+\beta_0(q_0')^2+\beta_\perp(Q')^2+\gamma_0 q_0'',\;\;
  \alpha\rho^2+\beta_\perp Q^2+\gamma_0\rho\,q_0',\;\;
  \alpha
\bigr).
\end{equation}
The Hessian contributes $\gamma_0 q_0''$ to $g_{\rho\rho}$ and
$\gamma_0\rho\,q_0'$ to $g_{\phi\phi}$.  Both are
$\rho$-dependent (whenever $q_0$ is non-constant), so the
metric is \emph{no longer flat}: its coefficients vary with
$\rho$, producing non-zero curvature.

\subsection{The modified Maxwell numerator}

\begin{theorem}[Maxwell numerator with Hessian generator]
\label{thm:hessian-numerator}
On the cylindrical ansatz~\eqref{eq:cyl-ansatz} with the
extended metric~\eqref{eq:cyl-metric-hess}, imposing the
wave equation $\square_g q_{2,3} = 0$ gives
\begin{equation}
\label{eq:hessian-numerator}
\nn_\mu F^{\mu\phi}
= \frac{\mathcal{N}}{g_{\rho\rho}\,g_{\phi\phi}^2},
\end{equation}
where
\begin{equation}
\label{eq:hessian-N}
\mathcal{N}
= \alpha\,(\rho Q'-Q)^2
  + \beta_0\,Q^2(q_0')^2
  + \gamma_0\,(\rho Q'-Q)\,(q_0'\,Q' - q_0''\,Q).
\end{equation}
\end{theorem}

\begin{proof}
The proof follows the product-rule strategy of
Theorem~\ref{thm:cyl-maxwell}.  The numerator is
$\mathcal{N} = Q^2 g_{\rho\rho} + (Q')^2 g_{\phi\phi}
- QQ'\,g_{\phi\phi}'$
with the extended metric components.  The $\beta_\perp$ terms
cancel exactly (as before).  The $\gamma_0$ contribution is:
\begin{align*}
&Q^2\,\gamma_0 q_0'' + (Q')^2\,\gamma_0\rho\,q_0'
- QQ'(\gamma_0 q_0' + \gamma_0\rho\,q_0'') \\
&\quad= \gamma_0\bigl[
  q_0''(Q^2 - \rho QQ') + q_0'(\rho(Q')^2 - QQ')
\bigr] \\
&\quad= \gamma_0\,(\rho Q'-Q)\,(q_0' Q' - q_0'' Q).
\end{align*}
\end{proof}

\begin{remark}[Breaking the sum-of-squares obstruction]
\label{rem:hessian-breaking}
In the order-$\le 1$ case (Theorem~\ref{thm:cyl-maxwell}),
the numerator $\mathcal{N} = \alpha(\rho Q'-Q)^2
+ \beta_0 Q^2(q_0')^2$ was a sum of non-negative terms,
forcing either $Q = c\rho$ (uniform field) or $q_0' = 0$
(no gravitational potential).  With the Hessian generator,
the cross-term $\gamma_0(\rho Q'-Q)(q_0' Q' - q_0'' Q)$
is \textbf{sign-indefinite}: it can be negative when
$\gamma_0$ and $(\rho Q'-Q)(q_0' Q' - q_0'' Q)$ have
opposite signs.  This breaks the no-go structure and opens
the door to solutions with \textbf{both} $Q \neq c\rho$
(non-uniform electromagnetic field) \textbf{and} $q_0' \neq 0$
(radial gravitational potential), producing a
\textbf{curved} self-consistent metric.
\end{remark}

\subsection{The coupled ODE system}

Setting $\mathcal{N} = 0$ and imposing the wave equations for
all four fields gives a coupled second-order ODE system for
$Q(\rho)$ and $q_0(\rho)$:
\begin{enumerate}[label=(\roman*)]
\item\label{item:wave-q0}
  \textbf{Wave equation for $q_0$:}
  $q_0' = q_0(\rho)$ depends only on $\rho$, so
  \begin{equation}
  \label{eq:wave-q0-hess}
  \frac{1}{\sqrt{|g|}}\,\pa_\rho\!\left(
    \frac{\sqrt{|g|}\,q_0'}{g_{\rho\rho}}\right) = 0,
  \qquad\text{i.e.}\qquad
  q_0'\,\sqrt{\frac{g_{\phi\phi}}{g_{\rho\rho}}} = C_0
  \end{equation}
  for a constant $C_0$.

\item\label{item:wave-Q}
  \textbf{Wave equation for $q_2,q_3$:}
  as in~\eqref{eq:wave-Q} (with $r\to\rho$),
  \begin{equation}
  \label{eq:wave-Q-hess}
  \frac{1}{\sqrt{|g|}}\,\pa_\rho\!\left(
    \frac{\sqrt{|g|}\,Q'}{g_{\rho\rho}}\right)
  = \frac{Q}{g_{\phi\phi}}.
  \end{equation}

\item\label{item:maxwell-hess}
  \textbf{Maxwell condition:}
  \begin{equation}
  \label{eq:maxwell-hess}
  \alpha\,(\rho Q'-Q)^2
  + \beta_0\,Q^2(q_0')^2
  + \gamma_0\,(\rho Q'-Q)(q_0'\,Q' - q_0''\,Q) = 0.
  \end{equation}
\end{enumerate}
The metric components
$g_{\rho\rho} = \alpha + \beta_0(q_0')^2 + \beta_\perp(Q')^2
+ \gamma_0 q_0''$ and
$g_{\phi\phi} = \alpha\rho^2 + \beta_\perp Q^2
+ \gamma_0\rho\,q_0'$
depend on the unknowns, making the system fully nonlinear and
self-referential---the metric is determined by the fields whose
dynamics it governs.  This is structurally identical to
the self-consistent system that produced curvature in the real
seam~\cite{seam1}.

\begin{remark}[The quantum-mechanical connection]
\label{rem:qm-connection}
The Hessian generator $\gamma_0\nn^2_h q_0$ makes the metric
respond to the \emph{curvature} of the scalar field---the
second-order structure that distinguishes wave-like solutions
from mere gradients.  In the Madelung representation of
quantum mechanics, the Hessian of the log-amplitude produces
the quantum potential; here, the Hessian of $q_0$ produces
gravitational curvature.  The wave equation
$\square_g q_0 = 0$ then closes the loop self-consistently:
the geometry curves because of the field's Hessian, and
the field's profile is shaped by propagation on that curved
geometry.  This self-referential coupling is the mechanism
that gave the real seam its non-trivial vacuum
solutions~\cite{seam1}, now extended to the quaternionic
setting where it simultaneously produces curvature
\emph{and} electromagnetism.
\end{remark}

\subsection{The curved cylindrical solution}
\label{sec:curved-cyl}

The coupled system~\eqref{eq:wave-q0-hess}--\eqref{eq:maxwell-hess}
is fully nonlinear, but the factored form of
the numerator~\eqref{eq:N-factored} suggests looking for
solutions where the bracket vanishes.  We now show that
this strategy yields an exact closed-form solution with
non-trivial curvature.

\begin{theorem}[Curved self-consistent Maxwell
  solution on a seam metric]
\label{thm:curved-cyl}
Setting $\beta_0 = 0$ in the extended $U(1)$-seam
metric~\eqref{eq:u1-seam-hessian} and choosing
$\beta_\perp c^2 \alpha = -\tfrac{1}{2}$,
$\gamma_0 = -\Lambda/d$ with $\Lambda > 0$, the fields
\begin{equation}
\label{eq:curved-fields}
q_0 = d\,\ln\rho, \qquad
q_1 = \omega\,t, \qquad
Q(\rho) = c\!\left(\alpha\rho
  - \frac{\Lambda}{\rho}\right),
\end{equation}
satisfy simultaneously:
\begin{enumerate}[noitemsep,label=(\roman*)]
\item the wave equations $\square_g q_a = 0$ for
  $a = 0,1,2,3$;
\item the source-free Maxwell equations
  $\nn_\mu F^{\mu\nu} = 0$ and $dF = 0$;
\item the seam metric is
\begin{equation}
\label{eq:curved-metric}
g = \operatorname{diag}\!\left(
  -\alpha+\beta_1\omega^2,\;\;
  \frac{\alpha^2\rho^4 - \Lambda^2}{2\alpha\rho^4},\;\;
  \frac{\alpha^2\rho^4 - \Lambda^2}{2\alpha\rho^2},\;\;
  \alpha
\right);
\end{equation}
\item the electromagnetic field is
\begin{equation}
\label{eq:curved-F}
F = \frac{c^2(\alpha^2\rho^4 - \Lambda^2)}{\rho^3}\,
  d\rho\wedge d\phi.
\end{equation}
\end{enumerate}
The metric degenerates at
$\rho_0 = (\Lambda/\alpha)^{1/2}$, where
$g_{\phi\phi} = g_{\rho\rho} = 0$ and the field strength
vanishes.  The Kretschner scalar diverges there
(Proposition~\ref{prop:curvature}), so $\rho_0$ is a
genuine curvature singularity---not a coordinate
artifact.  For $\rho > \rho_0$ the metric is Lorentzian
with $\rho$-dependent coefficients and non-zero
Riemann curvature.  The parameter $\Lambda$ controls
the strength of gravitational back-reaction: it
deforms the uniform vortex $Q = c\rho$ into a
non-uniform magnetic configuration whose core radius
$\rho_0$ grows with $\Lambda$.
\end{theorem}

\begin{proof}
\emph{Step~1.\ Maxwell condition.}
With $\beta_0 = 0$, the
numerator~\eqref{eq:hessian-N} factors:
\begin{equation}
\label{eq:N-factored}
\mathcal{N}
= (\rho Q'-Q)\bigl[\alpha(\rho Q'-Q)
  + \gamma_0(q_0' Q' - q_0'' Q)\bigr].
\end{equation}
The first factor $\rho Q'-Q = 2c\Lambda/\rho \neq 0$
(for the non-trivial branch $\Lambda > 0$), so we need
the second factor to vanish:
\begin{equation}
\label{eq:maxwell-bracket}
\alpha(\rho Q'-Q) + \gamma_0(q_0' Q' - q_0'' Q) = 0.
\end{equation}
With $q_0 = d\ln\rho$\,(so $q_0' = d/\rho$,
$q_0'' = -d/\rho^2$) and
$Q = c(\alpha\rho-\Lambda/\rho)$\,(so
$Q' = c(\alpha+\Lambda/\rho^2)$):
\begin{align*}
q_0' Q' - q_0'' Q
&= \frac{d}{\rho}\,c\!\left(\alpha+\frac{\Lambda}{\rho^2}\right)
  + \frac{d}{\rho^2}\,c\!\left(\alpha\rho
  - \frac{\Lambda}{\rho}\right) \\
&= \frac{cd}{\rho}\left[
  \alpha + \frac{\Lambda}{\rho^2}
  + \alpha - \frac{\Lambda}{\rho^2}\right]
= \frac{2cd\alpha}{\rho}.
\end{align*}
Substituting into~\eqref{eq:maxwell-bracket}:
\[
\alpha\cdot\frac{2c\Lambda}{\rho}
+ \gamma_0\cdot\frac{2cd\alpha}{\rho}
= \frac{2c\alpha}{\rho}\,(\Lambda + d\gamma_0) = 0.
\]
Since $\gamma_0 = -\Lambda/d$, this vanishes
identically.

\smallskip
\emph{Step~2.\ Wave equation for $q_0$.}
The extended metric components with $\beta_0 = 0$,
$\beta_\perp c^2 \alpha = -\tfrac{1}{2}$,
$\gamma_0 d = -\Lambda$ are:
\begin{align*}
g_{\rho\rho}
&= \alpha + \beta_\perp c^2\!\left(\alpha
  + \frac{\Lambda}{\rho^2}\right)^{\!2}
  + \gamma_0 q_0'' \\
&= \alpha - \frac{1}{2\alpha}\!\left(\alpha
  + \frac{\Lambda}{\rho^2}\right)^{\!2}
  + \frac{\Lambda}{\rho^2} \\
&= \frac{\alpha^2\rho^4 - \Lambda^2}{2\alpha\rho^4},
\end{align*}
and similarly
\[
g_{\phi\phi}
= \alpha\rho^2
  - \frac{1}{2\alpha}\!\left(\alpha\rho
  - \frac{\Lambda}{\rho}\right)^{\!2}
  - \Lambda
= \frac{\alpha^2\rho^4 - \Lambda^2}{2\alpha\rho^2}.
\]
Therefore $g_{\phi\phi}/g_{\rho\rho} = \rho^2$, and
the first integral~\eqref{eq:wave-q0-hess} gives
\[
q_0'\,\sqrt{\frac{g_{\phi\phi}}{g_{\rho\rho}}}
= \frac{d}{\rho}\cdot\rho = d = C_0. \qquad\checkmark
\]

\smallskip
\emph{Step~3.\ Wave equation for $Q$.}
With $g_{\phi\phi}/g_{\rho\rho} = \rho^2$,
the determinant factor
$\sqrt{|g|} = \sqrt{|g_{tt}|\,g_{\rho\rho}\,g_{\phi\phi}\,g_{zz}}$
has $g_{\rho\rho}\,g_{\phi\phi}\,g_{zz}$ proportional to
$(\alpha^2\rho^4 - \Lambda^2)^2/(4\alpha^2\rho^4)\cdot\alpha$,
so $\sqrt{|g|}\,g^{\rho\rho} \propto \rho$.
The wave equation~\eqref{eq:wave-Q-hess} reduces to
the Euler equation
\begin{equation}
\label{eq:euler-Q}
\rho^2\,Q'' + \rho\,Q' - Q = 0,
\end{equation}
whose general solution is $Q = c_1\rho + c_2/\rho$.
Our $Q = c(\alpha\rho - \Lambda/\rho)$ is precisely this,
with $c_1 = c\alpha$, $c_2 = -c\Lambda$.

\smallskip
\emph{Step~4.\ Non-trivial curvature.}
The Christoffel symbol
$\Gamma^\rho_{\phi\phi} = -g_{\phi\phi}'/(2g_{\rho\rho})
= -\rho(\alpha^2\rho^4 + \Lambda^2)/
(\alpha^2\rho^4 - \Lambda^2)
\neq -\rho$
(the flat-space value), so the Riemann tensor is non-zero
for $\Lambda \neq 0$. \qedhere
\end{proof}

\begin{figure}[t]
\centering
\includegraphics[width=\textwidth]{curved_metric.pdf}
\caption{Curved cylindrical solution
  (Theorem~\ref{thm:curved-cyl}) with
  $\alpha = \Lambda = c = 1$.
  \emph{Left:} metric components $g_{\rho\rho}$ (blue) and
  $g_{\phi\phi}$ (red), with the flat solution
  ($\Lambda = 0$, dotted) for comparison.  Both vanish at the
  critical radius $\rho_0 = \sqrt{\Lambda/\alpha}$ (dashed).
  \emph{Right:} field strength $F_{\rho\phi}$, showing the
  suppression near $\rho_0$ relative to the flat uniform field.}
\label{fig:curved-metric}
\end{figure}

\begin{remark}[Comparison with $Q = c\rho$]
\label{rem:curved-comparison}
Setting $\Lambda = 0$ gives $Q = c\alpha\rho$,
$g_{\phi\phi} = \alpha\rho^2/2$, $g_{\rho\rho} = \alpha/2$:
the flat solution of Proposition~\ref{prop:cyl-solution}
(up to a constant rescaling).  The parameter $\Lambda > 0$
deforms the flat vortex into a curved configuration
(Fig.~\ref{fig:curved-metric}):
\begin{itemize}[noitemsep]
\item The $1/\rho$ correction to $Q$ makes the magnetic
  field \emph{non-uniform}: $F_{\rho\phi}$ vanishes at the
  critical radius $\rho_0 = (\Lambda/\alpha)^{1/2}$ and
  reverses sign, akin to a magnetic null surface.
\item The metric develops a curvature singularity at
  $\rho_0$ where $g_{\phi\phi} = g_{\rho\rho} = 0$
  and the Kretschner scalar diverges
  (Proposition~\ref{prop:curvature}), forming a
  genuine geometric boundary.
\item The Hessian generator $\gamma_0\nn^2_h q_0$, with its
  logarithmic potential $q_0 = d\ln\rho$, is the mechanism
  that lifts $Q = c\rho$ to
  $Q = c(\alpha\rho - \Lambda/\rho)$.
\end{itemize}
\end{remark}

\begin{remark}[Sign of $\beta_\perp$]
\label{rem:beta-perp-sign}
The condition $\beta_\perp c^2 \alpha = -\tfrac{1}{2}$
requires $\beta_\perp < 0$.  This means the rotating-pair
gradient \emph{subtracts} from the metric, which is unusual
but physically natural: because the electromagnetic field is
part of the metric generation itself, a negative $\beta_\perp$
represents the \emph{binding energy} of the vortex---the
phase gradient pulls spacetime into a tight configuration
around the core, counteracting the positive conformal
factor~$\alpha$ that would otherwise make the space expand.
The seam coefficients are smooth functions of the jet
invariants, with no \emph{a~priori} sign constraint on
$\beta_\perp$.  The positivity of the metric is maintained
for $\rho > \rho_0 = (\Lambda/\alpha)^{1/2}$, where the
numerator $\alpha^2\rho^4 - \Lambda^2 > 0$ ensures both
$g_{\rho\rho} > 0$ and $g_{\phi\phi} > 0$.
Similarly, $\gamma_0 = -\Lambda/d < 0$, so the Hessian
generator subtracts from the metric.  Both signs are
consistent with the companion paper~\cite{seam1}, where
the Taub solution used either sign for the Hessian
coefficient~$\gamma$.
\end{remark}

% ------ Curvature analysis ------

\begin{proposition}[Curvature analysis of the curved solution]
\label{prop:curvature}
Consider the metric of Theorem~\ref{thm:curved-cyl} with
$H = \alpha^2\rho^4 - \Lambda^2$.
\begin{enumerate}[label=\textup{(\roman*)},noitemsep]
\item \textbf{Kretschner scalar.}
$K = R_{\mu\nu\rho\sigma}R^{\mu\nu\rho\sigma}
    = \frac{1024\,\Lambda^4\alpha^6\rho^{12}}{H^6}$.
Since $H \to 0$ as $\rho \to \rho_0$, we have
$K \to \infty$: the degeneration at $\rho_0$ is a
genuine curvature singularity.  Note that the
electromagnetic field $F_{\rho\phi} = c^2 H/\rho^3 \to 0$
vanishes at~$\rho_0$: the singularity is sourced by
the gravitational self-interaction of the metric
coefficients, not by concentration of
electromagnetic energy.

\item \textbf{Einstein tensor.}
$G_{\rho\rho} = G_{\phi\phi} = 0$,
while
$\displaystyle
  \frac{G_{tt}}{g_{tt}}
  = \frac{G_{zz}}{g_{zz}}
  = -\frac{16\,\Lambda^2\alpha^3\rho^6}{H^3}$.
Hence $G_{\mu\nu} = -\mathcal{K}(\rho)\,h^{(\|)}_{\mu\nu}$,
where $\mathcal{K} = 16\Lambda^2\alpha^3\rho^6/H^3 > 0$
is the Gaussian curvature of the $(\rho,\phi)$-surface
and $h^{(\|)}$ is the induced metric on the
$(t,z)$-plane orthogonal to the electromagnetic
field~$F = F_{\rho\phi}\,d\rho \wedge d\phi$.
In coordinate-invariant form,
$F_{\mu\alpha}\,G^{\alpha\nu} = 0$:
the Einstein tensor vanishes on every vector in
the support of~$F$.

\item \textbf{Relation to stress--energy.}
Standard Einstein--Maxwell requires
$G_{\mu\nu} = \kappa\,T^{(F)}_{\mu\nu}$ for some
$\kappa$.  Here $T^{(F)}_{\rho\rho}$ and
$T^{(F)}_{\phi\phi}$ are non-zero while
$G_{\rho\rho} = G_{\phi\phi} = 0$: no proportionality
constant satisfies the Einstein equation in all
components.  The same obstruction holds for the scalar
stress--energy $T^{(S)}_{\mu\nu}$ and for any linear
combination $T^{(F)} + T^{(S)}$.

This mismatch is \emph{structural}, not parametric:
the metric is a warped product
$g = g_{(t,z)} \oplus g_{(\rho,\phi)}$
with constant $(t,z)$-components, so all curvature
lives in the $(\rho,\phi)$-surface.  In two
dimensions the Einstein tensor vanishes
identically, forcing $G_{\rho\rho} = G_{\phi\phi} = 0$
regardless of the specific metric functions.  Since
$T^{(F)}$ is non-zero precisely in these directions,
no tuning of $\Lambda/\alpha$ can restore the
Einstein equation.

\item \textbf{Field-strength invariant.}
$F^{\alpha\beta}F_{\alpha\beta}
  = 8\alpha^2 c^4$,
independent of $\rho$---the EM field has constant
invariant magnitude despite living on a curved
background.  This contrasts with the Melvin magnetic
universe~\cite{Melvin1964}, where
$F^{\alpha\beta}F_{\alpha\beta}$ falls off as
$(1 + B_0^2\rho^2/4)^{-4}$.
\end{enumerate}
\end{proposition}

\begin{remark}[Self-consistency versus field equations]
\label{rem:no-einstein}
Item~(iii) shows that the seam framework does \emph{not}
reproduce the Einstein equation as an identity.
In general relativity, two independent postulates are
needed: (i)~matter tells spacetime how to curve
($G_{\mu\nu} = \kappa\,T_{\mu\nu}$), and
(ii)~spacetime tells matter how to move (field equations
on the curved background).  In the seam framework these
collapse into a single requirement: the self-consistency
loop---each scalar satisfies its wave equation on the
metric it generates---simultaneously determines both the
geometry and the field dynamics.
The emergent relation
$G_{\mu\nu} = -\mathcal{K}(\rho)\,h^{(\|)}_{\mu\nu}$
is neither trivial ($G = 0$) nor standard ($G = \kappa T$)
but exhibits \emph{geometric complementarity}: curvature
and electromagnetism occupy orthogonal subspaces of the
tangent bundle ($F_{\mu\alpha}\,G^{\alpha\nu} = 0$),
without requiring an extra dimension as in Kaluza--Klein
theory.
\end{remark}

% ================================================================
\section{Discussion}
\label{sec:discussion}
% ================================================================

The $U(1)$-equivariant quaternionic seam provides the natural
\emph{architecture} for coupling electromagnetism to
curvature:
$dF = 0$ is automatic, $T^{(F)}$ is traceless, and the
divergence equation reduces to an algebraic identity plus a
single geometric condition.  The self-consistency condition---that
each scalar field satisfies its wave equation on the metric it
generates---then acts as a \emph{symmetry filter}, selecting
the coordinate system compatible with emergent electromagnetism.

At order~$\le 1$ (gradient generators only), the cylindrical
vortex $Q = c\rho$ is the unique self-consistent Maxwell
solution, but it lives on a flat metric.  Restoring the
Hessian generator $\gamma_0\nn^2_h q_0$
(Section~\ref{sec:hessian}) breaks the sum-of-squares
obstruction in the Maxwell numerator.  The resulting coupled
ODE system admits an exact solution
(Theorem~\ref{thm:curved-cyl}):
$Q = c(\alpha\rho - \Lambda/\rho)$
with $q_0 = d\ln\rho$ produces a non-flat metric that
simultaneously satisfies the wave equations and source-free
Maxwell---a \emph{curved, self-consistent seam solution}
generated purely from scalars, in direct analogy with the
Melvin magnetic universe~\cite{Melvin1964}.  The curvature
analysis of Proposition~\ref{prop:curvature} shows that
the resulting Einstein tensor does \emph{not} satisfy
standard Einstein--Maxwell: $G_{\mu\nu}$ vanishes in the
electromagnetic plane while $T^{(F)}_{\mu\nu}$ does not.
Instead, $G_{\mu\nu}$ projects onto the $(t,z)$-subspace
orthogonal to~$F$, a geometric relation imposed by
self-consistency rather than postulated.

\begin{center}
\resizebox{\textwidth}{!}{%
\begin{tabular}{lccl}
\textbf{Ansatz} & \textbf{Background} & \textbf{Maxwell?}
  & \textbf{Obstruction / solution} \\
\hline
$Q(r)(\cos\phi,\sin\phi)$
  & spherical & \textbf{No}
  & $\cos^2\!\theta$ mismatch (Thm~\ref{thm:no-go-maxwell}) \\
$Q(\rho)(\cos\phi,\sin\phi)$
  & cylindrical & \textbf{Yes}
  & $Q = c\rho$, uniform $\mathbf{B}$, flat
    (Thm~\ref{thm:cyl-maxwell}) \\
$Q(r)(\cos\theta,\sin\theta)$
  & spherical & \textbf{No}
  & eigenvalue $\ell(\ell+1)=2$ too strong
    (Thm~\ref{thm:polar-nogo}) \\
Hedgehog $f(r)\hat{x}_a$
  & spherical & ---
  & requires non-abelian (Rem.~\ref{rem:hedgehog}) \\
Cylindrical $+$ Hessian
  & cylindrical & \textbf{Yes}
  & $Q = c(\alpha\rho-\Lambda/\rho)$, curved
    (Thm~\ref{thm:curved-cyl}) \\
\end{tabular}%
}
\end{center}

\medskip\noindent
\textbf{Open questions.}

\begin{enumerate}[noitemsep]
\item \textbf{Monopole from intermediate equivariance.}
  A spherically symmetric monopole requires non-abelian field
  strength.  Does an intermediate equivariance group (between
  $U(1)$ and $Sp(1)$) support a self-consistent monopole
  within seam geometry?

\item \textbf{Non-flat backgrounds.}
  All results use a flat Minkowski background $h$.  Can the
  seam framework produce non-flat vacuum solutions (e.g.\ Kerr)
  or self-consistent configurations on curved backgrounds?

\item \textbf{Modified field equation.}
  Proposition~\ref{prop:curvature} shows $G_{\mu\nu} \neq
  \kappa\,T^{(F)}_{\mu\nu}$ for any $\kappa$: the seam
  framework does not reproduce the standard Einstein equation
  as an identity.  The emergent relation
  $G_{\mu\nu} = -\mathcal{K}(\rho)\,h^{(\|)}_{\mu\nu}$, or
  equivalently $F_{\mu\alpha}\,G^{\alpha\nu} = 0$, projects the
  Einstein tensor onto the plane orthogonal to~$F$.
  Can this projection identity be derived from an action
  principle, and does it generalise beyond cylindrical
  symmetry?

\item \textbf{Massive fields and quantum potential.}
  This paper uses massless wave equations $\square_g q_a = 0$.
  The companion paper~\cite{seam1} showed that the massive
  Klein--Gordon system introduces a nonlinear
  amplitude--frequency coupling; in the quaternionic setting,
  replacing the massless equations with a quaternionic
  Klein--Gordon system would couple the electromagnetic and
  gravitational sectors through the mass term, potentially
  linking photon dispersion to spacetime curvature.
\end{enumerate}

% ================================================================
\backmatter

\begin{thebibliography}{99}

\bibitem{seam1}
R\"onnb\"ack, L.: The Einstein tensor in seam geometry: generator
closure, emergent signature, and null-convexity.
Submitted to Gen. Relativ. Gravit. (2026)

\bibitem{BaezHuerta2010}
Baez, J.C., Huerta, J.: Division algebras and supersymmetry~I.
 In: Doran, R., Friedman, G., Rosenberg, J. (eds.) Superstrings,
 Geometry, Topology, and $C^*$-Algebras, Proc.\ Sympos.\ Pure
 Math., vol.~81, pp.~65--80. Amer.\ Math.\ Soc., Providence (2010)

\bibitem{Kaluza1921}
Kaluza, T.: Zum Unit\"atsproblem der Physik.
Sitzungsber.\ Preuss.\ Akad.\ Wiss.\ Berlin (Math.\ Phys.)
\textbf{1921}, 966--972 (1921)

\bibitem{Klein1926}
Klein, O.: Quantentheorie und f\"unfdimensionale
Relativit\"atstheorie.
Z.\ Phys.\ \textbf{37}(12), 895--906 (1926)

\bibitem{Rainich1925}
Rainich, G.Y.: Electrodynamics in the general relativity theory.
Trans.\ Amer.\ Math.\ Soc.\ \textbf{27}(1), 106--136 (1925)

\bibitem{MisnerWheeler1957}
Misner, C.W., Wheeler, J.A.: Classical physics as geometry:
Gravitation, electromagnetism, unquantized charge, and mass as
properties of curved empty space.
Ann.\ Phys.\ \textbf{2}(6), 525--603 (1957)

\bibitem{tHooft1974}
't~Hooft, G.: Magnetic monopoles in unified gauge theories.
Nucl.\ Phys.\ B \textbf{79}(2), 276--284 (1974)

\bibitem{Polyakov1974}
Polyakov, A.M.: Particle spectrum in quantum field theory.
JETP Lett.\ \textbf{20}(6), 194--195 (1974)

\bibitem{Melvin1964}
Melvin, M.A.: Pure magnetic and electric geons.
Phys.\ Lett.\ \textbf{8}(1), 65--68 (1964)

\end{thebibliography}

\end{document}
